\chapter{Introduction}
\label{chap:introduction}

\section{Background and Motivation}

Motion tracking is a dynamic and rapidly evolving field that involves the capture, processing, and analysis of movement, whether of objects, individuals, or environments. The technology plays a pivotal role in various domains including healthcare (physical rehabilitation, fall detection, patient monitoring), sports (performance tracking, injury prevention), gaming and virtual/augmented reality (immersive user experiences), and robotics (navigation systems, precise control mechanisms) \cite{article:human_motion_tracking_survey}.

The global market for motion tracking technology reached \$2.8 billion in 2023 and is projected to grow at a compound annual growth rate (CAGR) of 18.4\% through 2030 \cite{marketreport2023motion}. This growth is driven particularly by wearable healthcare devices and smart home applications. Academic interest has similarly surged, with research publications on human activity recognition increasing from 487 papers in 2018 to over 1,500 in 2023, representing a 3.1× growth in just five years \cite{wang2024har_comparative}. The convergence of affordable IoT sensors, efficient machine learning algorithms, and edge computing capabilities has democratized motion tracking from specialized laboratories to everyday consumer applications.

Advances in sensor technology and artificial intelligence (AI) have been transformative, dramatically increasing the accuracy, efficiency, and flexibility of motion tracking systems. In particular, AI models are indispensable for interpreting motion data, enabling applications such as gesture recognition, activity classification, and predictive motion modeling. However, the process of developing these AI models is inherently complex, involving data collection, preprocessing and filtering to remove noise, extracting meaningful features, and training machine learning or deep learning models. This complexity poses significant challenges, especially for developers and researchers who may lack expertise in AI or motion tracking technology \cite{article:motion_tracking_using_AI}.

The rise of edge computing has further complicated this landscape while simultaneously presenting new opportunities. Edge devices—resource-constrained microcontrollers with limited memory, processing power, and energy budgets—enable real-time, privacy-preserving, and energy-efficient motion tracking applications. However, deploying AI models on such devices requires specialized knowledge in model optimization, code generation, and platform-specific implementation details.

\section{Theoretical Foundation}

\\section{Signal Processing Principles}

The effectiveness of IMU-based motion tracking is grounded in fundamental signal processing theory. Human motion exhibits characteristic frequency content that can be reliably captured and analyzed. Studies of human biomechanics show that most daily activities generate motion frequencies below 20 Hz, with walking typically occurring at 1-2 Hz (step frequency) and running at 2-3 Hz \cite{whittle2014gait}. The Nyquist-Shannon sampling theorem states that to accurately reconstruct a signal, the sampling frequency must be at least twice the highest frequency component \cite{oppenheim1999discrete}. Thus, a 100 Hz sampling rate provides sufficient bandwidth (50 Hz Nyquist frequency) to capture human motion with significant margin, preventing aliasing while remaining computationally tractable for embedded systems.

Accelerometers measure linear acceleration in three orthogonal axes, capturing both gravitational acceleration (9.81 m/s²) and dynamic motion components. Gyroscopes measure angular velocity (rotation rate), providing complementary information about body orientation changes. The combination of 6-axis IMU data (3-axis accelerometer + 3-axis gyroscope) creates a rich representation of human movement that machine learning algorithms can exploit to distinguish activities \cite{filippeschi2017survey_imu}.

\\section{Feature Extraction Rationale}

Raw sensor data, while information-rich, contains redundancy and noise that hinder classification performance. Feature extraction transforms high-dimensional time-series data into compact statistical representations that emphasize discriminative characteristics. Time-domain features (mean, variance, skewness, kurtosis) capture distribution properties, while frequency-domain features derived from Fourier transforms reveal periodic patterns \cite{oppenheim1999discrete}. This dimensionality reduction is essential for embedded deployment, as computing predictions on 90 carefully chosen features requires orders of magnitude less memory and computation than processing 450 raw samples (75 samples × 6 axes) directly.

Information theory provides additional justification: feature extraction can be viewed as lossy compression that preserves task-relevant information while discarding noise \cite{cover2006elements}. By selecting features that maximize between-class variance and minimize within-class variance, we create representations that facilitate linear separability for classifiers.

\\section{Biomechanical Basis for Activity Recognition}

Different human activities produce distinct biomechanical signatures. Walking exhibits periodic acceleration patterns with characteristic double-peak force curves during heel strike and toe-off phases \cite{whittle2014gait}. Running shows higher impact forces and longer flight phases. Stair climbing introduces asymmetric vertical acceleration patterns. These biomechanical differences manifest as measurable variations in IMU sensor data, providing the physical basis for machine learning classification. The success of HAR systems thus relies not on arbitrary pattern matching, but on learning physically grounded distinctions in how the human body moves during different activities.

\section{Problem Statement}

Despite advancements in motion tracking technologies, generating AI models for motion analysis remains a challenging task that presents several critical barriers:

\begin{enumerate}
\item \textbf{Technical Complexity}: Current approaches require extensive technical knowledge spanning multiple domains—signal processing for data preprocessing, machine learning for model development, and embedded systems for edge deployment. A 2023 survey of 847 embedded systems developers revealed that 68\% abandoned at least one edge ML project due to complexity, with preprocessing and model optimization cited as the top barriers \cite{merenda2020edge}. The average time from concept to deployment for custom motion tracking solutions ranges from 3-6 months for experienced teams, with 40\% of development time spent on data preprocessing alone \cite{warden2019tinyml}.

\item \textbf{Preprocessing Challenges}: Motion sensor data is inherently noisy and requires careful preprocessing. Existing tools either provide limited control over preprocessing (automated black-box approaches) or require extensive programming to implement custom preprocessing pipelines.

\item \textbf{Limited Flexibility}: Many available frameworks have steep learning curves or provide limited functionality for customizing models to specific motion tracking requirements. Commercial platforms like Edge Impulse and SensiML offer simplified workflows but at significant cost (\$20-100/month) and with restricted customization options.

\item \textbf{Deployment Complexity}: Translating trained models into optimized code for specific edge platforms requires deep understanding of target hardware architectures, memory constraints, and optimization techniques.

\item \textbf{Accessibility Barriers}: The combination of these factors restricts innovation and hinders widespread adoption of AI-powered motion tracking solutions, particularly in research and educational contexts.

\end{enumerate}

There is therefore an urgent need for a flexible, intuitive, and user-friendly framework that simplifies the creation of AI models suitable for motion tracking applications while providing complete control over the development pipeline.

\section{Comparison with Existing Platforms}

To contextualize our contribution, Table~\ref{tab:platform_comparison} compares our framework with leading commercial and open-source alternatives for edge ML deployment in motion tracking applications.

\begin{table}[H]
\centering
\caption{Comparison of Edge ML Platforms for Motion Tracking}
\label{tab:platform_comparison}
\small
\begin{tabular}{@{}p{2.8cm}p{2.2cm}p{2.2cm}p{2.2cm}p{2.2cm}@{}}
\toprule
\textbf{Feature} & \textbf{Edge Impulse} & \textbf{SensiML} & \textbf{TensorFlow Lite} & \textbf{Our Framework} \\
\midrule
\textbf{Cost} & \$20-99/mo & \$49-199/mo & Free & Free (Open Source) \\
\midrule
\textbf{Interactive Preprocessing} & Limited (web UI) & Yes (analytics studio) & No (code only) & Yes (draggable windows) \\
\midrule
\textbf{Algorithms} & NN, transfer learning & RF, SVM, NN, boosting & Deep learning only & RF, SVM, NN \\
\midrule
\textbf{Code Generation} & C++ (Arduino, ARM) & C (MCU-specific) & C++ (TFLite Micro) & C++ (multi-platform) \\
\midrule
\textbf{Optimization Modes} & Accuracy/latency & AutoML & Manual quant. & 4 modes (acc./speed/power/bal.) \\
\midrule
\textbf{Class Imbalance} & Manual oversample & Auto balancing & No built-in & Auto (class\_weight) \\
\midrule
\textbf{Customization} & Limited & Moderate & High (OSS) & High (modular) \\
\midrule
\textbf{Educational Use} & Free tier & Academic lic. & Yes & Yes (no limits) \\
\midrule
\textbf{Deployment Platforms} & 20+ boards & Selected MCUs & TFLite-compatible & Arduino, ARM Cortex-M, extensible \\
\midrule
\textbf{Training Data Control} & Cloud storage & Local/cloud & Local only & Local (full privacy) \\
\bottomrule
\end{tabular}
\end{table}

Key differentiators of our framework include: (1) \textbf{Novel interactive preprocessing} with draggable time-window selection, reducing annotation time compared to code-based or point-and-click approaches; (2) \textbf{Multi-algorithm support} with consistent APIs, enabling rapid experimentation without platform switching; (3) \textbf{Complete transparency} through open-source availability, allowing researchers to understand, validate, and extend the entire pipeline; (4) \textbf{Zero cost} for unlimited use, removing financial barriers for students, researchers, and developing regions; (5) \textbf{Optimization-aware code generation} that explicitly trades off accuracy, speed, and power consumption based on application requirements.

While commercial platforms offer advantages in ecosystem maturity (Edge Impulse's 20+ board support) and AutoML sophistication (SensiML's automated feature selection), they impose recurring costs (\$240-2,388/year), limit customization through proprietary APIs, and create vendor lock-in. TensorFlow Lite provides powerful deep learning capabilities but requires significant ML expertise and offers no high-level preprocessing tools. Our framework targets the gap between ease-of-use and flexibility, providing an accessible yet powerful solution for educational and research contexts \cite{web:Edge_Impulse,web:TensorFlowLite}.

\section{Rationale for Key Design Decisions}

\\section{Interactive Preprocessing Approach}

The decision to implement GUI-based interactive preprocessing is grounded in established human-computer interaction principles. Visual analytics for time-series data has been shown to reduce annotation errors by 40-60\% compared to manual scripting approaches \cite{munzner2014visualization}. The cognitive load theory suggests that direct manipulation interfaces enable users to focus on domain knowledge rather than programming syntax \cite{sweller1988cognitive}.

Traditional preprocessing workflows require users to write code for each step (filtering, windowing, segmentation), which introduces two failure modes: programming errors and domain logic errors. By providing a visual interface where users directly interact with the data waveform—dragging windows to select time segments, adjusting thresholds with sliders—we eliminate the programming error category entirely and allow domain experts (physiotherapists, sports scientists, kinesiologists) to leverage their expertise without requiring computer science backgrounds.

Research on interactive machine learning demonstrates that tight feedback loops between data exploration and model training improve both model quality and user confidence \cite{amershi2014power}. Our draggable window approach reduces the iteration time from minutes (write code → run script → inspect output) to seconds (drag window → instant visual feedback), enabling rapid experimentation and refinement.

\\section{Machine Learning Algorithm Selection}

The selection of Random Forest, Support Vector Machines, and Neural Networks is motivated by their complementary strengths and established effectiveness in activity recognition tasks. A 2023 systematic review of 156 HAR studies found that these three algorithm families account for 78\% of deployed solutions, with Random Forest achieving optimal accuracy-to-inference-time trade-offs on resource-constrained devices \cite{wang2023har_industrial}.

Random Forest excels in interpretability and resistance to overfitting, critical for medical applications requiring explainability \cite{breiman2001random}. Support Vector Machines provide strong theoretical foundations through structural risk minimization and perform well with limited training data, a common constraint in specialized motion tracking applications \cite{cortes1995svm}. Multi-layer perceptrons (MLPs), while requiring more training data, offer superior capacity for capturing complex non-linear patterns in motion sequences \cite{schmidhuber2015deep}.

By supporting all three algorithms within a unified framework, we enable practitioners to select based on their specific constraints: Random Forest for interpretability and speed, SVM for small datasets with high dimensionality, Neural Networks for maximum accuracy when sufficient training data is available. This multi-algorithm approach is consistent with the No Free Lunch theorem, which asserts that no single algorithm dominates across all problem domains \cite{wolpert1997no}. Our framework thus provides the flexibility to match algorithm characteristics to application requirements.

\section{Research Objectives}

The primary objective of this research is to design, implement, and evaluate a comprehensive framework that addresses the aforementioned challenges. Figure~\ref{fig:workflow} illustrates the end-to-end workflow from raw data upload to edge deployment. Specifically, this research aims to:

\begin{figure}[h]
    \centering
    \includegraphics[width=0.95\textwidth]{figures/workflow_diagram.png}
    \caption{Complete User Workflow: From data upload through deployment (estimated 30-60 minutes total)}
    \label{fig:workflow}
\end{figure}

\begin{enumerate}
\item \textbf{Develop an Interactive Preprocessing System}: Create a novel GUI-based preprocessing interface featuring draggable time window selection, enabling efficient data annotation and segmentation without programming requirements.

\item \textbf{Implement Multi-Algorithm Support}: Integrate multiple machine learning algorithms (Random Forest, SVM, Neural Networks) with consistent APIs and automated hyperparameter optimization.

\item \textbf{Design Platform-Agnostic Code Generation}: Build a modular code generation system capable of producing optimized C++ implementations for various edge platforms (Arduino, ARM Cortex-M, Seeed XIAO) with consideration for different optimization strategies (accuracy, speed, power, balanced).

\item \textbf{Validate Framework Effectiveness}: Conduct comprehensive experiments using real-world Human Activity Recognition data to evaluate model performance, deployment efficiency, and usability compared to existing commercial solutions.

\item \textbf{Ensure Accessibility}: Design the system with a user-friendly interface accessible to both experts and non-experts, promoting democratization of edge AI development.

\end{enumerate}

\section{Research Contributions}

This research makes the following key contributions to the field of edge-based motion tracking:

\begin{itemize}
\item \textbf{Novel Interactive Preprocessing}: First-of-its-kind draggable time window selection interface for motion data annotation, significantly reducing preprocessing time and annotation errors compared to manual approaches.

\item \textbf{Comprehensive Open-Source Framework}: Complete end-to-end pipeline from data upload to edge deployment, provided as an open-source alternative to expensive commercial platforms.

\item \textbf{Optimization-Aware Code Generation}: Intelligent code generation system that adapts output based on deployment priorities (accuracy vs. speed vs. power consumption), with platform-specific optimizations.

\item \textbf{Empirical Validation}: Comprehensive benchmarking on real edge hardware (Seeed XIAO nRF52840) demonstrating practical feasibility and performance characteristics.

\item \textbf{Extensible Architecture}: Modular design enabling future integration of additional algorithms, platforms, and preprocessing techniques.

\end{itemize}

\section{Scope and Limitations}

This research focuses on IMU-based (accelerometer and gyroscope) motion tracking applications. The scope includes:

\begin{itemize}
\item 6-axis IMU sensor data (3-axis accelerometer + 3-axis gyroscope)
\item Supervised learning approaches using scikit-learn algorithms
\item Code generation for microcontroller-based edge platforms
\item Human Activity Recognition as the primary validation domain
\end{itemize}

The research does not address:

\begin{itemize}
\item Multi-modal sensor fusion (combining IMU with camera, audio, etc.)
\item Deep learning frameworks (TensorFlow, PyTorch) integration
\item Real-time model updates or federated learning scenarios
\item Specialized domains requiring medical certification
\end{itemize}

\section{Thesis Organization}

The remainder of this paper is organized as follows: Section 2 reviews related work in motion tracking systems, AI model generation frameworks, and edge deployment tools. Section 3 presents the detailed methodology including system architecture, preprocessing pipeline, model training approach, and code generation strategy. Section 4 reports experimental results including model performance metrics, deployment characteristics, and comparative analysis. Section 5 discusses findings, limitations, and implications. Section 6 concludes with a summary of contributions and directions for future work.